\scanpage{29}
\mspage{3}
\uline{Corollaire} 2 Soit $\deleted{\unsure{X|A}}A$ un schéma abélien sur $S$ loc. noethérien, et $G$ un préschéma en groupes loc. de type fini sur $S$. Alors \deleted{\unsure{tout}} \deleted{\unsure{$S$-}}hom \inserted{de $S$-préschémas} de $A$ dans $G$ est de la forme (de façon unique)
\[
f=u+c
\]
\margin{Attention, si $G$ n'est pas commutatif}
% slanted
où $u$ est un hom en groupes, et où $c$ est un hom ``constant'' \inserted{i.e. provient d'une section de $G$ sur $S$}. On a nécessairement $c=f\circ\varepsilon_A$, où $\varepsilon_A$ est la section unité de $A$. \unsure{Donc} si $f\circ\varepsilon_A=\varepsilon_G$, $f$ est un hom de schémas en groupes.

\uline{Corollaire} 3 Soit $A$ un schéma abélien sur $S$. Alors $A$ est commutatif. Toute loi de composition sur $A$, telle que $\varepsilon_A$ \deleted{\unsure{soit}} soit unité bilatère, est la loi de composition donnée de $A$.

Il suffit de prouver cette dernière assertion. Or on a, si $f:A\times A\longrightarrow A$ est la loi de composition, en vertu du cor. 1 :
\[
f=[f\circ(\varepsilon\times\operatorname{id}_A)]p_1+
(f\circ(\operatorname{id}_A\times\varepsilon))p_2+
\deleted{\unsure{f\circ(\varepsilon_1}}[f\circ(\varepsilon_1\times\varepsilon_2)]p_2,
\]
i.e. $f=p_1+p_2$, c'est la loi de composition donnée. OK,
% archive sheet number: 28
